% Options for packages loaded elsewhere
\PassOptionsToPackage{unicode}{hyperref}
\PassOptionsToPackage{hyphens}{url}
\documentclass[
  ignorenonframetext,
]{beamer}
\newif\ifbibliography
\usepackage{pgfpages}
\setbeamertemplate{caption}[numbered]
\setbeamertemplate{caption label separator}{: }
\setbeamercolor{caption name}{fg=normal text.fg}
\beamertemplatenavigationsymbolsempty
% remove section numbering
\setbeamertemplate{part page}{
  \centering
  \begin{beamercolorbox}[sep=16pt,center]{part title}
    \usebeamerfont{part title}\insertpart\par
  \end{beamercolorbox}
}
\setbeamertemplate{section page}{
  \centering
  \begin{beamercolorbox}[sep=12pt,center]{section title}
    \usebeamerfont{section title}\insertsection\par
  \end{beamercolorbox}
}
\setbeamertemplate{subsection page}{
  \centering
  \begin{beamercolorbox}[sep=8pt,center]{subsection title}
    \usebeamerfont{subsection title}\insertsubsection\par
  \end{beamercolorbox}
}
% Prevent slide breaks in the middle of a paragraph
\widowpenalties 1 10000
\raggedbottom
\AtBeginPart{
  \frame{\partpage}
}
\AtBeginSection{
  \ifbibliography
  \else
    \frame{\sectionpage}
  \fi
}
\AtBeginSubsection{
  \frame{\subsectionpage}
}
\usepackage{iftex}
\ifPDFTeX
  \usepackage[T1]{fontenc}
  \usepackage[utf8]{inputenc}
  \usepackage{textcomp} % provide euro and other symbols
\else % if luatex or xetex
  \usepackage{unicode-math} % this also loads fontspec
  \defaultfontfeatures{Scale=MatchLowercase}
  \defaultfontfeatures[\rmfamily]{Ligatures=TeX,Scale=1}
\fi
\usepackage{lmodern}
\ifPDFTeX\else
  % xetex/luatex font selection
\fi
% Use upquote if available, for straight quotes in verbatim environments
\IfFileExists{upquote.sty}{\usepackage{upquote}}{}
\IfFileExists{microtype.sty}{% use microtype if available
  \usepackage[]{microtype}
  \UseMicrotypeSet[protrusion]{basicmath} % disable protrusion for tt fonts
}{}
\makeatletter
\@ifundefined{KOMAClassName}{% if non-KOMA class
  \IfFileExists{parskip.sty}{%
    \usepackage{parskip}
  }{% else
    \setlength{\parindent}{0pt}
    \setlength{\parskip}{6pt plus 2pt minus 1pt}}
}{% if KOMA class
  \KOMAoptions{parskip=half}}
\makeatother
\setlength{\emergencystretch}{3em} % prevent overfull lines
\providecommand{\tightlist}{%
  \setlength{\itemsep}{0pt}\setlength{\parskip}{0pt}}
\usepackage{bookmark}
\IfFileExists{xurl.sty}{\usepackage{xurl}}{} % add URL line breaks if available
\urlstyle{same}
\hypersetup{
  hidelinks,
  pdfcreator={LaTeX via pandoc}}

\author{\texorpdfstring{}{}}
\date{}

\begin{document}

\begin{frame}{Abstract}
\protect\phantomsection\label{sec:abstract}
Active Inference provides a unified formalism for understanding agents
that minimize variational free energy through perception and action.
While most treatments focus on Active Inference as a first-order theory
of surprise minimization, comparatively little attention has been paid
to its operation at the \emph{meta-level}---as a methodology that is
simultaneously \emph{meta-pragmatic} (specifying what matters) and
\emph{meta-epistemic} (specifying what can be known). This gap limits
our understanding of how framework specification itself shapes cognitive
outcomes and constrains the space of possible inferences.

This paper addresses this gap by introducing a \(2 \times 2\) matrix
framework that organizes Active Inference's contributions along two
orthogonal dimensions: Data versus Meta-Data and Cognitive versus
Meta-Cognitive processing. The resulting four quadrants---baseline EFE
computation (Q1), meta-data-weighted processing (Q2), reflective
self-monitoring (Q3), and framework-level optimization (Q4)---provide a
systematic taxonomy for analyzing how Active Inference operates across
multiple levels of abstraction.

We demonstrate that the Expected Free Energy (EFE) formulation operates
at a meta-level where modeler choices in specifying matrices \(A\),
\(B\), \(C\), and \(D\) define the boundaries of both epistemic and
pragmatic domains. Unlike reinforcement learning, where reward functions
are externally imposed, Active Inference makes framework specification
itself a research variable. Computational validation through
quadrant-specific implementations confirms that meta-data integration
improves decision reliability (85\% to 94\%), meta-cognitive control
enhances robustness under uncertainty, and framework-level optimization
yields a 23\% overall performance improvement.

The implications extend to cognitive security, where meta-level
processing reveals novel attack surfaces and defense strategies across
each quadrant, and to AI safety, where principled framework
specification provides a foundation for value alignment that transcends
scalar reward functions. By foregrounding the meta-level structure of
Active Inference, this work opens new avenues for understanding
intelligence as framework flexibility---the capacity to modify the very
structures within which cognition operates.

\textbf{Keywords:} active inference, free energy principle,
meta-cognition, meta-pragmatic, meta-epistemic, cognitive science,
cognitive security, framework specification, generative models

\textbf{MSC2020:} 68T01 (Artificial intelligence), 91E10 (Cognitive
science), 92B05 (Neural networks)
\end{frame}

\end{document}
